\title{Residual Ternary Algebra and a modular foundation:\\ the $s$-series, the Bridge Series $C$, and the 35-component law vector}
\author{Arvin Hampton\\ \small 539 Labs}
\date{27 October 2026 (v14)}
\begin{document}
\maketitle
\begin{abstract}
Residual Ternary Algebra (RTA, also called Resonant Algebra) is built from three kinds of written objects: elementary packaging arithmetic
(the ternary map $T_3$, the residual cardinality $539$, the carrier $\mathcal O_{\rm res}$), the residual series
$g_s(q)=\sum_{n\ge0}q^{n^2}/(-q^s;q^s)_n$, and the Bridge Series $C=\Porth_7-H_7$, an orthant-minus-cone sum on the indefinite lattice
$\Z^2$ with Gram matrix $\left(\begin{smallmatrix}2&7\\7&7\end{smallmatrix}\right)$. This paper collects, with complete proofs, what can be
established in tying these objects to a modular foundation, and what cannot.

Part~I treats the $s$-series $g_s$ (written $g_k$ there): an exact expansion of the coefficients, families of zeros, sign laws for even $k$,
and, for every odd $k\le33$, a parity-dominance theorem with complete zero sets (computer-assisted in stated steps). Part~II treats the Bridge
Series $C$: a ray-class decomposition over $\Q(\sqrt{35})$, radial asymptotics, a Baker-type obstruction to every scalar modular or mock residual law,
and a vector-valued structure from a Siegel theta lift. Part~III is new. It proves that the leading radial profile of $C$ lies outside the span
$\VV$ of the 35-component law vector, while the profile of the normalized series $q^{-7/20}C$ lies in $\VV$ and obeys its vector-valued
$S$-law; it gives the complete radial expansion of $q^{-7/20}C$ and of $105$ companion series through order $t^0$ at every rational point,
with closed formulas, an algebraicity statement at order $t^{-1/2}$, and an exact decomposition of that coefficient by quadratic class.
The final part states precisely which links between RTA and this framework are proved, which written RTA objects ($T_3$, $\mathcal O_{\rm res}$,
$539$, the packaging identities) have no written link to it, and where the shadow $\Xi$ would have to enter; the last question is open.
Part~III also determines the shadow structure of the $105$ companion series: the exact matrix identities $M^2=35P$ and $(MT)^3=35\sqrt{35}P$, the resulting parity
theorem (a weight-one law forces the trivial multiplier and $P$-odd vectors), the $\xi$-images of the cone completions in terms of three unary theta series, the Laplacian
defect of the orthant Siegel lift as one product of unary thetas with an exact vanishing criterion, and the line terms of the literal series in closed form. It shows that the transformation law of the cone completions is exactly the law of the $105$ companions, that the completed odd part has trivial radial expansion at every rational point, and that the order $t^{-1/2}$ coefficient equals the sum of explicit completion and line terms; at order $t^0$ the coefficient is the sum of line terms and an explicit even corner term, and no affine law of weight-one type holds for the latter; its Bernoulli-weighted parts obey an exact $S$-law into Gauss means outside the companion family. The corner term is a cone invariant, transported exactly by the automorphs of the lattice; the unit acts on the companions as an involution commuting with the $S$-matrix, and the orbit of the corner under the full automorph group generated by the unit and a reflection has two elements; the residual of the fixed $S$-matrix law for real-weight sums, including $A_2$, satisfies an exact $x\leftrightarrow1/x$ symmetry.
This version records the October 2026 packaging identities for the local map $A_1\to\Xi\to\mathrm{RTA}$ (Section~\ref{sec:partial13}): outer, path/branch/$\psi$, and $\varepsilon\to0^+$/W7 identities, together with non-local items A1--A12. The Bridge link remains partial: no readiness claim is made and the residual question for $C$ remains open. The $h$-table of item G-A3 is paused with an empty schema, item X4 is absent with a stated gap, and no written samples of $h_{\rm residual}^{\rm high}$ exist; these are catalogue-only entries, not theorems.
This version further develops, in Section~\ref{sec:modlocal}, the local modular constituents that complete the modular side of the foundation: equality of compact and ordinary induction for characters inflated through the torus; reduction of the $K$-socle to $\Gamma$-socles of $K_1$-invariants; the Breuil--M\'ezard / Colmez dictionary identifying $B(1,\omega)$ with $\pi_\alpha$ and $V(\beta)$ with $\rho$; the determination $w(\theta^3 f)=10$ of the Serre--Katz filtration weight; Breuil's socle formula for $\pi(r,0,\chi)$ at every $p\ge3$; and the classification of irreducible $\overline{\mathbb{F}}_p$-representations of $\mathrm{GL}_2(\mathbb{F}_p)$ as $\mathrm{Sym}^s\otimes\det^a$, obtained without Brauer's theorem. These identities do not settle the residual question for~$C$ and do not identify the shadow~$\Xi$. No claim is made about the twin prime conjecture or the Riemann hypothesis.
\end{abstract}
\tableofcontents
\clearpage
\section*{Symbol key}\addcontentsline{toc}{section}{Symbol key}
\begin{center}\small
\begin{longtable}{p{3.2cm}p{9.3cm}p{2.4cm}}\toprule
symbol & meaning & where defined\\\midrule\endhead
\multicolumn{3}{l}{\emph{General and preliminaries}}\\
$q=\e(\tau)$, $\e(z)$ & $\e(z)=e^{2\pi iz}$ & \S\ref{sec:prelim}\\
$Q$, $B$, $A$ & binary quadratic form $n_1^2+7n_1n_2+\frac72n_2^2$, its bilinear form, Gram matrix & \S\ref{sec:prelim}\\
$Q_s$ & quadratic form of parameter $s$ & Def.~\ref{def:series}\\
$\Porth_s$ & orthant theta series of $Q_s$ & Def.~\ref{def:series}\\
$\Pi_s$ & product series $(-q^s;q^s)_\infty g_s$ & Prop.~\ref{prop:PiPorth}\\
$H_s$ & cone series with walls $c_1,c_2$ (never redefined) & \S\ref{sec:conv}\\
$C$ & Bridge Series $\Porth_7-H_7$ & \S\ref{sec:p3lead}\\
$w(v)$ & cone weight $\frac12(\sgn B(c_1,v)-\sgn B(c_2,v))$ & \S\ref{sec:p3lead}\\
$I$ & $\log(6+\sqrt{35})/(2\sqrt{35})$ & Thm.~\ref{thm:lead}\\
\multicolumn{3}{l}{\emph{Part I}}\\
$g_k$, $g_k(d)$ & residual series $\sum q^{n^2}/(-q^k;q^k)_n$, its coefficients & \S\ref{definitions}\\
$P_n$, $\varphi_n(m)$ & $\prod_{j\le n}(1+x^j)^{-1}$ and its coefficients & \S\ref{definitions}\\
$A_n(m)$, $Q(m)$ & $|\varphi_n(m)|$; partitions of $m$ into distinct odd parts & \S\ref{the-polynomials-varphi_n}\\
$Z_k$ & zero set of $g_k$ on allowed $d$ & \S\ref{definitions}\\
$\mathcal S_k$, $\mathrm{OccFilt}_k(\infty)$ & square classes mod $k$; classes eventually nonzero & \S\ref{definitions}\\
$\Xi_k$ & a finite set of integers (unrelated to the shadow $\Xi$) & \S\ref{the-cases-k15-and-k21}\\
\multicolumn{3}{l}{\emph{Part II}}\\
$g_s$ & residual series at parameter $s$ & \S\ref{app:gs}\\
$V$ & $42$-dim.\ span of the lift functions & \S\ref{sec:lift}\\
$E$ & ray-class partial theta series & \S\ref{sec:rq}\\
$\Caxis$, $\Cthird$ & pieces of $C$ & \S\ref{sec:cthird}\\
$u_L$ & vector of the theta-lift family & \S\ref{sec:uL}\\
$\Delta$ & wall parameter of the cone & \S\ref{sec:delta}\\
$H_s^{\rm axes}$ & axis variant of $H_s$ & \S\ref{sec:axes}\\
\multicolumn{3}{l}{\emph{Part III}}\\
$G_{\alpha,\beta}$ & Gauss mean & Def.~\ref{def:gaussmean}\\
$D$ (set), $a$, $h_A$ & discriminant group; shift; $(\frac12,0)$ & \S\ref{sec:law35}\\
$K1_\mu,\dots,K4_\mu$ & $G_{a+\mu,a},G_{a+\mu,0},G_{\mu,a},G_{\mu,0}$ & Def.~\ref{def:K}\\
$(K2_\mu)$ & the 35-component law vector & Def.~\ref{def:K}\\
$K2_0$, $K3_p$ & components at $\mu=0$, $\mu=p$ & Def.~\ref{def:K}\\
$\VV$ & $36$-dim.\ span of $K2,K3,K4$ & Def.~\ref{def:K}, Prop.~\ref{prop:dimV}\\
$\EE$, $\CC$ & $36$- and $6$-dim.\ summands of Part II's $42$-dim.\ lift rep. & Prop.~\ref{prop:sl2}\\
$S_{\alpha,\beta}$, $\wt$ & companion series and weight $1_{\rm orthant}-w$ & Def.~\ref{def:companion}\\
$q^{-7/20}C=S_{a,0}$ & normalized Bridge Series & \S\ref{sec:p3lead}\\
$u_1,u_2$; $K_\pm$ & $(21,-1)$, $(0,1)$; cones split into $21$ pieces & Prop.~\ref{prop:decomp}\\
$\Lambda_{\alpha,\beta}$ & leading radial profile & Thm.~\ref{thm:p3profile}\\
$\Ar$, $J_a$, $J_b$, $K(l)$ & Gaussian cone integrals & Lem.~\ref{lem:integrals}\\
$\mean(f)$, $W_i$, $V_i$, $\mathsf C$ & periodic mean, sawtooth, second moment, corner terms & Thm.~\ref{thm:piece}\\
$A_0,A_1,A_2$ & coefficients of $t^{-1},t^{-1/2},t^0$ & Thm.~\ref{thm:expansion}\\
$D$ (term), $R$, $W[\lambda]$ & drift, ray, sawtooth parts of $A_1$ & Prop.~\ref{prop:A1split}\\
$R[\frac72]$, $R[\frac{595}2]$ & ray terms on $\pm\N(0,1)$, $\pm\N(21,-1)$ & Prop.~\ref{prop:A1split}\\
$\lambda\in\{1,\frac72,\frac{595}2\}$ & $Q$-values $Q(e_1)$, $Q(0,1)$, $Q(21,-1)$ & Prop.~\ref{prop:A1split}\\
$Z_0$ & zero set of $A_0$ in $X_{216}$ & Prop.~\ref{prop:A1nonzero}\\
$\Xi_\varepsilon$ & sawtooth shadow candidate (Category B) & \S\ref{sec:order4}\\
\multicolumn{3}{l}{\emph{Local modular constituents (Part~III)}}\\
$G$, $B=P$ & $\mathrm{GL}_2(\mathbb{Q}_p)$; upper-triangular Borel & \S\ref{sec:modlocal}\\
$\chi$, $T$ & smooth character of $P$ inflated through the torus $T$ & \S\ref{sec:modlocal}\\
$\mathrm{Ind}_B^G$, $\mathrm{c\text{-}Ind}_B^G$ & ordinary and compact induction & Thm.~\ref{thm:cind-ind}\\
$K$, $K_1$, $\Gamma$ & $\mathrm{GL}_2(\mathcal{O}_F)$; first congruence kernel; $\mathrm{GL}_2(\mathbb{F}_p)$ & Thm.~\ref{thm:soc-k1}\\
$\mathrm{soc}_K$, $\mathrm{soc}_\Gamma$ & socle of a smooth representation & Thm.~\ref{thm:soc-k1}\\
$\pi_\alpha$ & $\mathrm{Ind}_P^G\,\omega\otimes\omega^{-1}$ (Breuil--M\'ezard) & Thm.~\ref{thm:colmez-d12}\\
$B(\delta_1,\delta_2)$, $V(\beta)$ & Colmez principal series; nonsplit extension attached to $\beta$ & Thm.~\ref{thm:colmez-d12}\\
$\omega$, $\rho$ & mod-$p$ cyclotomic character; nonsplit extension $0\to\omega\to\rho\to1\to0$ & Thm.~\ref{thm:colmez-d12}\\
$\theta$, $w(\theta^3 f)$ & Serre--Katz theta operator; filtration weight of $\theta^3 f$ & Prop.~\ref{prop:filtr-bounds}, Thm.~\ref{thm:filtr-10}\\
$\pi(r,0,\chi)$ & Breuil principal series $(\mathrm{ind}_{KZ}^G\mathrm{Sym}^r/T)\otimes\chi\circ\det$ & Thm.~\ref{thm:breuil-soc}\\
$\mathrm{Sym}^s\otimes\det^a$ & irreducible of $\mathrm{GL}_2(\mathbb{F}_p)$ ($0\le s\le p-1$) & Thm.~\ref{thm:irreps-gl2}\\
$I(1)$ & pro-$p$ Iwahori subgroup of $K$ & Prop.~\ref{prop:irreps-k}\\
\multicolumn{3}{l}{\emph{Part IV}}\\
$T_3$ & ternary map $3k\mapsto k$, $3k{+}1\mapsto4k{+}2$, $3k{+}2\mapsto2k{+}1$ & Prop.~\ref{prop:T3}\\
$\mathcal O_{\rm res}$ & carrier $\{21+9j:0\le j\le538\}$ & Rem.~\ref{rem:nolink}\\
$539$ & residual cardinality & Prop.~\ref{prop:T3}\\
$\Xi$ & RTA weight-$\frac32$ shadow (not identified) & \S\ref{sec:found}\\
$M,T,P$ & $105\times105$ $S$-, $T$- and parity matrices on the companions & \S\ref{sec:shadow}\\
$p,E,\rho_O,\rho_H$ & $\nu$-sign weight, its even part, odd cone weights & Lemma~\ref{lem:psplit}\\
$\widehat H,\widehat G_O,X^c$ & cone completions and their $\xi_1$-images & Prop.~\ref{prop:xi}\\
$\mathcal B_{z,z'},\vartheta^\pm$ & boundary series; unary thetas in $\frac72m^2$, $\frac52u^2$ & Thm.~\ref{thm:Bc1}\\
$\Phi_{\rm orth}$ & Siegel lift over the orthant arc & Thm.~\ref{thm:lap}\\
$\bar B_1,\kappa_c,R_c$ & periodic sawtooth; wall weights; wall corrections & \S\ref{sec:radial}\\
$A_2^{\rm lines},A_2^{\rm corner}$ & order-$t^0$ line and corner terms & Thm.~\ref{thm:A2exact}\\
$A_1^{\rm lines}$ & order-$t^{-1/2}$ line contribution & Thm.~\ref{thm:A1exact}\\
$\rho_m,\rho_u$ & reflections $m\mapsto-m$, $u\mapsto-u$ & Thm.~\ref{thm:Bc1}\\
\bottomrule
\end{longtable}
\end{center}
\clearpage

